\documentclass[a4paper]{article}
\usepackage{amsfonts}
\usepackage{amsmath}
\usepackage{amssymb}
\usepackage{graphicx}     

\begin{document}
  
\noindent \large Auteur: Odia Kibor \\
\noindent \large Studierichting: Informatica\\
\noindent \large Oefeningenreeks 2C 8.17 \\

\medskip

\normalsize

$2^{3x} + 8^{x} = 27^{x} + 3^{3x} + \left(\sqrt{3}\right)^{6x}$ \\

$\Leftrightarrow 2^{3x} + \left(2^{3}\right)^{x} = \left(3^{3}\right)^{x} + 3^{3x} + \left(3^{\frac{1}{2}}\right)^{6x}$ \\

$\Leftrightarrow 2^{3x} + 2^{3x} = 3^{3x} + 3^{3x} + 3^{3x}$ \\

$\Leftrightarrow 2 \cdot 2^{3x} = 3 \cdot 3^{3x}$ \\

$\Leftrightarrow 2^{1+3x} = 3^{1+3x}$ \\

Deze twee termen kunnen enkel gelijk zijn als de exponent gleijk is aan $0$ \\

$\Rightarrow 1 + 3x = 0$ \\

$\Rightarrow x = \dfrac{-1}{3}$ \\


\begin{thebibliography}{999}
%\bibitem{a} Referentie
\end{thebibliography}
\end{document} 
